% !TeX root = ../beamer.tex

\section{Maths and code}

\subsection{Theorem}
\begin{frame}{Theorem}
    \begin{theorem}[Pythagoras]
        $a^{2}+b^{2}=c^{2}$
    \end{theorem}

    \begin{theorem}[A formula for the stress test]
        \[
            \mathcal{F}[f](\xi)=\int_{-\infty}^{\infty} f(x)\,e^{-2\pi i x\xi}\,dx,\qquad
            \zeta(s)=\sum_{n=1}^{\infty}\frac{1}{n^{s}}=\prod_{p\ \mathrm{prime}}\frac{1}{1-p^{-s}}
        \]
        \[
            \hat{\boldsymbol\beta}=\bigl(\mathbf{X}^{\top}\mathbf{X}\bigr)^{-1}\mathbf{X}^{\top}\mathbf{y},\quad
            \mathbf{A}=\begin{pmatrix}a_{11}&\cdots&a_{1n}\\\vdots&\ddots&\vdots\\a_{m1}&\cdots&a_{mn}\end{pmatrix},\quad
            \binom{n}{k}=\frac{n!}{k!\,(n-k)!}
        \]
        \[
            \lim_{n\to\infty}\Pr\!\Bigl(\frac{\bar{X}_n-\mu}{\sigma/\sqrt{n}}\le z\Bigr)
            =\frac{1}{\sqrt{2\pi}}\int_{-\infty}^{z}e^{-t^{2}/2}\,dt=\Phi(z)
        \]
    \end{theorem}
\end{frame}

\subsection{Definitions and examples}
\begin{frame}{Definitions and examples}
\begin{columns}[T,onlytextwidth]
\begin{column}{0.5\textwidth}
\step{1. A definition}
\begin{codeonly}
\begin{definition}[Title]
  Text with \alert{emphasis}
\end{definition}
\end{codeonly}
\end{column}
\begin{column}{0.45\textwidth}
\stepoffset
\begin{definition}[Estimator]
    A rule that turns a sample into a guess is an \alert{estimator}.
\end{definition}
\end{column}
\end{columns}
\bigskip
\begin{columns}[T,onlytextwidth]
\begin{column}{0.5\textwidth}
\step{2. An example}
\begin{codeonly}
\begin{example}
  Text
\end{example}
\end{codeonly}
\end{column}
\begin{column}{0.45\textwidth}
\stepoffset
\begin{example}
    The sample mean $\bar{x}$ estimates $\mu$.
\end{example}
\end{column}
\end{columns}
\end{frame}

\subsection{Theorems and proofs}
\begin{frame}{Theorems and proofs}
\begin{columns}[T,onlytextwidth]
\begin{column}{0.5\textwidth}
\step{1. A lemma, with its proof}
\begin{codeonly}
\begin{lemma}
  Statement
  \begin{proof}
    Proof text
  \end{proof}
\end{lemma}
\end{codeonly}
\end{column}
\begin{column}{0.45\textwidth}
\stepoffset
\begin{lemma}
    $\operatorname{Var}(aX)=a^{2}\operatorname{Var}(X)$.
    \begin{proof}
        Expand and take $a^{2}$ out.
    \end{proof}
\end{lemma}
\end{column}
\end{columns}
\bigskip
\begin{columns}[T,onlytextwidth]
\begin{column}{0.5\textwidth}
\step{2. A theorem or a corollary}
\begin{codeonly}
\begin{theorem}[Title]
\begin{corollary}
\end{codeonly}
\end{column}
\begin{column}{0.45\textwidth}
\stepoffset
\begin{theorem}[Pythagoras]
    $a^{2}+b^{2}=c^{2}$.
\end{theorem}
\end{column}
\end{columns}
\end{frame}

\subsection{Without a box}
\begin{frame}{Without a box}
\begin{columns}[T,onlytextwidth]
\begin{column}{0.5\textwidth}
\step{1. Add a star}
\begin{codedoc}
\begin{theorem*}[Pythagoras]
  $a^{2}+b^{2}=c^{2}$.
  \begin{proof}
    Compare the areas.
  \end{proof}
\end{theorem*}
\end{codedoc}
\end{column}
\begin{column}{0.45\textwidth}
\stepoffset
\begin{theorem*}[Pythagoras]
    $a^{2}+b^{2}=c^{2}$.
    \begin{proof}
        Compare the areas.
    \end{proof}
\end{theorem*}
\end{column}
\end{columns}
\bigskip
\begin{columns}[T,onlytextwidth]
\begin{column}{0.5\textwidth}
\step{2. Also definitions and examples}
\begin{codedoc}
\begin{definition*}[Estimator]
\begin{example*}
\end{codedoc}
Lemmas and corollaries too.
\end{column}
\begin{column}{0.45\textwidth}
\stepoffset
\begin{definition*}[Estimator]
    A rule that turns a sample into a guess.
\end{definition*}
\begin{example*}
    The sample mean $\bar{x}$ estimates $\mu$.
\end{example*}
\end{column}
\end{columns}
\end{frame}

\subsection{Formulas: code and result}
\begin{frame}{Formulas: code and result}
\begin{columns}[T,onlytextwidth]
\begin{column}{0.5\textwidth}
\step{1. A theorem}
\begin{codeonly}
\begin{theorem}[Title]
  $\bar{X}_n \xrightarrow{p} \mu$
\end{theorem}
\end{codeonly}
\end{column}
\begin{column}{0.45\textwidth}
\stepoffset
\begin{theorem}[Large numbers]
    $\bar{X}_n \xrightarrow{p} \mu$
\end{theorem}
\end{column}
\end{columns}
\bigskip
\begin{columns}[T,onlytextwidth]
\begin{column}{0.5\textwidth}
\step{2. A formula}
\begin{codeonly}
\[
  t=\frac{\bar{x}-\mu_0}{s/\sqrt{n}}
\]
\end{codeonly}
\end{column}
\begin{column}{0.45\textwidth}
\stepoffset
\[
  t=\frac{\bar{x}-\mu_0}{s/\sqrt{n}}
\]
\end{column}
\end{columns}
\end{frame}

\subsection{Code boxes}
\begin{frame}{Code boxes}
\begin{columns}[T,onlytextwidth]
\begin{column}{0.5\textwidth}
\step{1. Show code without running it}
\begin{codedoc}
\begin{codeonly}
  \alert{some code}
\end{codeonly}
\end{codedoc}
\end{column}
\begin{column}{0.45\textwidth}
\stepoffset
\begin{codeonly}
\alert{some code}
\end{codeonly}
\end{column}
\end{columns}
\bigskip
\begin{columns}[T,onlytextwidth]
\begin{column}{0.5\textwidth}
\step{2. Code in a sentence}
\begin{codedoc}
The function \code{mean(xs)} returns
the average, and in Python
\code[Python]{return sum(xs)}.
\end{codedoc}
\end{column}
\begin{column}{0.45\textwidth}
\stepoffset
The function \code{mean(xs)} returns the average, and in Python \code[Python]{return sum(xs)}.
\end{column}
\end{columns}
\end{frame}

\subsection{Source code}
\begin{frame}{Source code}
\begin{columns}[T,onlytextwidth]
\begin{column}{0.5\textwidth}
\step{1. A block per language}
\begin{codedoc}
\begin{codeblock}{Python}
  def mean(xs):
      return sum(xs) / len(xs)
\end{codeblock}
\end{codedoc}
\smallskip
\step{2. Any language of listings}
C, C++, Python and Julia are set up. For another one, name it as listings does.
\end{column}
\begin{column}{0.45\textwidth}
\stepoffset
\begin{codeblock}{Python}
# the average of a list
def mean(xs):
    return sum(xs) / len(xs)

print(mean([1, 2, 3]))  # 2.0
\end{codeblock}
\end{column}
\end{columns}
\end{frame}

\subsection{Pseudocode}
\begin{frame}{Pseudocode}
\begin{columns}[T,onlytextwidth]
\begin{column}{0.5\textwidth}
\step{1. In the font of the deck}
\begin{codedoc}
\begin{pseudocode}
  m <- A[1]
  for each a in A do
    if a > m then m <- a
\end{pseudocode}
\end{codedoc}
Keywords are bold, \code{<-} is an arrow, and a formula goes between dollar signs.
\end{column}
\begin{column}{0.45\textwidth}
\stepoffset
\begin{pseudocode}
function max_of(A)
    m <- A[1]            // best so far
    for each a in A do
        if a > m then m <- a
    return m
\end{pseudocode}
\end{column}
\end{columns}
\bigskip
\begin{columns}[T,onlytextwidth]
\begin{column}{0.5\textwidth}
\step{2. In typewriter, as a code box}
\begin{codedoc}
\begin{codeblock}{Pseudocode}
  while lo <= hi do
    mid <- $\lfloor (lo+hi)/2 \rfloor$
\end{codeblock}
\end{codedoc}
\end{column}
\begin{column}{0.45\textwidth}
\stepoffset
\begin{codeblock}{Pseudocode}
while lo <= hi do
    mid <- $\lfloor (lo+hi)/2 \rfloor$
\end{codeblock}
\end{column}
\end{columns}
\end{frame}
